\section{Trace-pairing expansion in transfer-matrix form}
\label{sec:trace_pairing_expansion}
%% =========================================================

% Formula/source: Wolf, Eq. (2.20), writes the transfer coefficients for
% Hilbert--Schmidt orthonormal bases as
% \widehat T_{\alpha\beta}=\tr(F_\alpha^\dagger T(G_\beta)).  Identifying the
% two bases with a Hermitian Hilbert--Schmidt orthonormal basis {\sigma_i}
% gives t_{ij}=\tr(\sigma_i T(\sigma_j)).  The chapter proves the analogous
% coordinate expansion more generally for a basis self-dual under the
% bilinear pairing (A,B)\mapsto\tr(AB):
% T(\rho)=\sum_{i,j}t_{ij}\tr(\sigma_j\rho)\sigma_i.
% Index verdict: j is summed in the input coordinate \tr(\sigma_j\rho), while
% i is summed in the coefficient and output basis element \sigma_i.  This is an
% algebraic basis expansion, not one tensor-network contraction.  In the
% retired figure only \sigma_j and \rho were contracted; t_{ij}, \sigma_i, and
% the equality to T(\rho) were unattached labels.

\begin{definition}[Trace-self-dual basis]\label{def:trace_pairing_self_dual_basis}
    \lean{TracePairingSelfDualBasis}
    \leanok
    A basis $\{\sigma_i\}_i$ of $\MN{D}$ is \emph{trace-self-dual} when
    its coordinate functionals are given by trace pairing:
    \[
        X = \sum_i \tr(\sigma_i X) \sigma_i.
    \]
\end{definition}

\begin{definition}[Trace-pairing coefficients]\label{def:trace_pairing_coeff}
    \lean{tracePairingCoeff}
    \leanok
    \uses{def:trace_pairing_self_dual_basis}
    For a linear map $T : \MN{D} \to \MN{D}$ and a trace-self-dual basis
    $\{\sigma_i\}$, define coefficients
    \[
        t_{ij} := \tr(\sigma_i\,T(\sigma_j)).
    \]
\end{definition}

\begin{theorem}[Trace-pairing expansion of a linear map]\label{thm:linearMap_expand_tracePairing}
    \lean{linearMap_expand_tracePairing}
    \leanok
    \uses{def:trace_pairing_self_dual_basis, def:trace_pairing_coeff}
    If $\{\sigma_i\}$ is trace-self-dual, then every linear map
    $T : \MN{D} \to \MN{D}$ admits the expansion
    \[
        T(\rho)
        = \sum_i \sum_j t_{ij}
        \tr(\sigma_j \rho) \sigma_i,
        \qquad
        t_{ij} = \tr(\sigma_i T(\sigma_j)).
    \]
\end{theorem}

\begin{proof}
    \leanok
    Expand $\rho$ and each $T(\sigma_j)$ in the chosen basis and substitute.
    The trace-self-dual identity replaces coordinates by traces, giving the
    stated double-sum formula.
\end{proof}

%% =========================================================
\section{SVD normal form (existence)}
\label{sec:svd_normal_form}
%% =========================================================

\cite[Section~2.3]{Wolf2012Quantum} discusses two families of normal forms for
the transfer-matrix representation of a quantum channel: the \emph{singular
    value decomposition} (SVD) and, for qubit channels, the \emph{Lorentz normal
    form} obtained by general invertible Kraus-rank-one CP filtering operations.
The SVD normal form is proved below. The Lorentz normal form theorem
(Theorem~\ref{thm:lorentz_normal_form_qubit}) uses general invertible
Kraus-rank-one completely positive filterings. Its proof requires the minimisation
argument and Lorentz-orbit analysis in
\cite[Propositions~2.9 and~2.11]{Wolf2012Quantum}; see
Section~\ref{sec:lorentz_normal_form}.

\begin{theorem}[SVD for positive semidefinite matrices]
    \label{thm:svd_posSemidef}
    \lean{Matrix.svd_of_posSemidef}\leanok
    Every positive semidefinite matrix $M \in \MN{D}$ admits a decomposition
    \[
        M = U \diag(\sigma) U^{\dagger},
    \]
    where $U \in \mc{U}(D)$ is unitary and $\sigma : \{0, \ldots, D-1\} \to
        \R_{\ge 0}$ has non-negative entries.
\end{theorem}

\begin{proof}\leanok
    The spectral theorem for Hermitian matrices provides an orthonormal
    eigenbasis with real eigenvalues. Positive semidefiniteness forces those
    eigenvalues to be non-negative, and the decomposition is written in SVD
    form with $U$ the eigenvector unitary and $\sigma$ the eigenvalue sequence.
\end{proof}

\begin{theorem}[SVD existence for invertible matrices]
    \label{thm:svd_of_isUnit}
    \lean{Matrix.svd_of_isUnit}\leanok
    Every invertible complex square matrix $M \in \MN{D}$ admits a singular
    value decomposition
    \[
        M = U \diag(\sigma) V^{\dagger},
    \]
    with $U, V \in \mc{U}(D)$ unitary and $\sigma_i > 0$ for all $i$.
\end{theorem}

\begin{proof}\leanok
    Apply the spectral theorem to the positive definite matrix
    $M^{\dagger} M$ to obtain $M^{\dagger} M = V \diag(\lambda)
        V^{\dagger}$ with $\lambda_i > 0$. Set $\sigma_i := \sqrt{\lambda_i}$ and
    $\Sigma := \diag(\sigma)$; since $\Sigma$ is a real positive
    diagonal, it is self-adjoint and invertible. Define $U := M V \Sigma^{-1}$.
    Then $U^{\dagger} U = \Sigma^{-1} V^{\dagger} (M^{\dagger} M) V \Sigma^{-1}
        = \Sigma^{-1} \diag(\lambda) \Sigma^{-1} = \Id$, so $U$ is
    unitary, and $U \Sigma V^{\dagger} = M V (\Sigma^{-1} \Sigma) V^{\dagger} =
        M V V^{\dagger} = M$.
\end{proof}

\begin{theorem}[SVD representation of a transfer matrix]
    \label{thm:transferMatrix_svd_of_isUnit}
    \lean{transferMatrix_svd_of_isUnit}\leanok
    \uses{thm:svd_of_isUnit}
    Every invertible transfer matrix $\widehat{T}$ of a linear super-operator
    $T : \MN{D} \to \MN{D}$ admits a singular value decomposition
    $\widehat{T} = U \diag(\sigma) V^{\dagger}$ with $U, V$
    unitary on $\C^{D} \otimes \C^{D}$ and $\sigma_i > 0$.
\end{theorem}

\begin{proof}\leanok
    Direct specialisation of Theorem~\ref{thm:svd_of_isUnit} to the index type
    $\{0,\ldots,D-1\} \times \{0,\ldots,D-1\}$ used by the transfer-matrix
    representation.
\end{proof}

\begin{remark}[Sorted / unique singular values]
    Theorem~\ref{thm:svd_of_isUnit} produces the singular values as an
    unordered family, without the usual convention that $\sigma_{i}$ are sorted
    in non-increasing order or the statement that they are uniquely determined
    by $M$. Downstream applications in \cite[Section~2.3]{Wolf2012Quantum},
    such as trace-norm identities, polar decomposition, and the Lorentz normal
    form, will typically require the sorted / unique variant; that refinement is
    future work.
\end{remark}

%% =========================================================
\section{Lorentz normal form}
\label{sec:lorentz_normal_form}
%% =========================================================

This section states the existence theorems from
\cite[Section~2.3]{Wolf2012Quantum}.
% The existing Wolf.* declarations referenced in this section live in
% TNLean/Channel/LorentzNormalForm.lean. The correctly formulated Proposition 2.11
% has no Lean declaration yet; declaration tags are omitted for pending results.

\begin{definition}[SL-filtering operation]\label{def:sl_filtering}
    \lean{Wolf.SLFiltering}\leanok
    An \emph{SL-filtering} for $D \times D$ matrices is a completely
    positive map of the form $\Phi(X) = S X S^{\dagger}$ where
    $S \in \MN{D}$ satisfies $\det S = 1$. Such maps are
    invertible and have Kraus rank~1.
\end{definition}

\begin{definition}[Doubly-stochastic map]\label{def:doubly_stochastic}
    \lean{Wolf.DoublyStochastic}\leanok
    A linear map $T : \MN{D} \to \MN{D}$ is \emph{doubly-stochastic}
    if $T(\Id) \propto \Id$ and the reduced density matrix
    $\tr_{1}[\tau]$ of its Choi matrix $\tau = (T \otimes \id)(\ketbra{\Omega}{\Omega})$
    is proportional to the identity. This is the normal form in
    \cite[Proposition~2.9]{Wolf2012Quantum}.
\end{definition}

\begin{lemma}[Positive-definite trace lower bound]\label{lem:posdef_trace_lower_bound}
    \lean{posDef_minEigenvalue_mul_trace_conjTranspose_mul_self_le,
        sub_minEigenvalue_smul_one_posSemidef}
    \leanok
    Let $D\geq 1$, let $M\in \MN{D}$ be positive-definite, and let
    $\lambda_{\min}(M)$ be its smallest Hermitian eigenvalue. For every
    $X\in \MN{D}$,
    \[
        \lambda_{\min}(M)\,\tr(X^{\dagger}X)
        \leq
        \tr(XMX^{\dagger}).
    \]
\end{lemma}

\begin{proof}
    \leanok
    The matrix $M-\lambda_{\min}(M)\Id$ is positive semidefinite by the
    spectral theorem. Since $X^{\dagger}X$ is also positive semidefinite, the
    trace product of these two matrices is non-negative. Expanding
    \[
        \tr\!\left(X^{\dagger}X(M-\lambda_{\min}(M)\Id)\right)
    \]
    and cycling the trace gives the displayed inequality.
\end{proof}

\begin{lemma}[Determinant AM--GM Hilbert--Schmidt bound]
\label{lem:det_amgm_hilbert_schmidt_bound}
    \lean{Matrix.card_le_trace_conjTranspose_mul_self_re_of_det_norm_eq_one}
    \leanok
    Let $A\in \MN{D}$ with $D\geq 1$. If $|\det A|=1$, then
    \[
        D \leq \tr(A^{\dagger}A).
    \]
\end{lemma}

\begin{proof}
    \leanok
    The matrix $A^{\dagger}A$ is positive semidefinite and has determinant
    $|\det A|^{2}=1$. Thus the product of its Hermitian eigenvalues is one.
    Denoting the eigenvalues of $A^{\dagger}A$ by
    $\lambda_{1},\ldots,\lambda_{D}\geq 0$, the arithmetic-geometric mean
    inequality gives
    \[
        \frac{\lambda_{1}+\cdots+\lambda_{D}}{D}
        \geq(\lambda_{1}\cdots\lambda_{D})^{1/D}=1,
    \]
    so $\tr(A^{\dagger}A)=\lambda_{1}+\cdots+\lambda_{D}\geq D$.
\end{proof}

\begin{lemma}[Kronecker Hilbert--Schmidt trace multiplicativity]
\label{lem:kronecker_hilbert_schmidt_trace}
    \lean{Matrix.trace_conjTranspose_mul_self_kronecker}
    \leanok
    For complex matrices $A$ and $B$ of compatible rectangular sizes,
    \[
        \tr\!\left((A\otimes_{k}B)^{\dagger}(A\otimes_{k}B)\right)
        =
        \tr(A^{\dagger}A)\,\tr(B^{\dagger}B).
    \]
\end{lemma}

\begin{proof}
    \leanok
    Use
    $(A\otimes_{k}B)^{\dagger}=A^{\dagger}\otimes_{k}B^{\dagger}$, the mixed
    product identity for Kronecker products, and
    $\tr(C\otimes_{k}D)=\tr(C)\tr(D)$. With
    $C=A^{\dagger}A$ and $D=B^{\dagger}B$, these identities give the equality.
\end{proof}

% Lean: Wolf.infimum_is_attained
\begin{lemma}[Infimum attainment for SL-filterings]\label{lem:infimum_is_attained}
    \lean{Wolf.infimum_is_attained}\leanok
    \uses{lem:posdef_trace_lower_bound, lem:det_amgm_hilbert_schmidt_bound,
    lem:kronecker_hilbert_schmidt_trace}
    For a positive-definite Choi matrix $\tau$, the infimum of
    $\tr\![(S_{2} \otimes_{k} S_{1}) \tau (S_{2} \otimes_{k} S_{1})^{\dagger}]$
    over $S_{1}, S_{2} \in \MN{D}$ with $\det S_{1} = \det S_{2} = 1$
    is attained.
\end{lemma}
\begin{proof}\leanok
    \uses{lem:posdef_trace_lower_bound, lem:det_amgm_hilbert_schmidt_bound,
    lem:kronecker_hilbert_schmidt_trace}
    Write $X = S_{2} \otimes_{k} S_{1}$. By the positive-definite trace lower
    bound (Lemma~\ref{lem:posdef_trace_lower_bound}),
    $\tr[X \tau X^{\dagger}] \ge \lambda_{\min}(\tau)\,\tr[X^{\dagger}X]$, and
    by Hilbert--Schmidt multiplicativity
    (Lemma~\ref{lem:kronecker_hilbert_schmidt_trace}),
    $\tr[X^{\dagger}X] = \tr[S_{2}^{\dagger}S_{2}]\,\tr[S_{1}^{\dagger}S_{1}]$.
    The determinant AM--GM estimate
    (Lemma~\ref{lem:det_amgm_hilbert_schmidt_bound}) gives
    $\tr[S_{i}^{\dagger}S_{i}] = \|S_{i}\|^{2} \ge D$ whenever
    $\det S_{i} = 1$. Combining, the value dominates
    $\lambda_{\min}(\tau)\,D\,\|S_{i}\|^{2}$ for each factor, so any point
    whose value is at most the value at the identity lies in a fixed
    Frobenius ball $\{\,\|S\| \le C\,\}$. Intersecting with the closed set
    $\det S = 1$ produces a compact set, on which the continuous trace
    functional attains its minimum by the extreme-value theorem. A point
    outside the sublevel set has value larger than the value at the identity,
    so the minimiser on the compact set is a global minimiser.
\end{proof}

% Lean: pow_card_mul_prod_le_sum_pow
\begin{lemma}[Arithmetic--geometric-mean inequality, product/sum form]
\label{lem:amgm_prod_sum}
    \lean{pow_card_mul_prod_le_sum_pow}
    \leanok
    For a nonnegative family $f_{0},\dots,f_{D-1}$ of real numbers,
    \[
        D^{D}\prod_{i} f_{i} \le \Bigl(\sum_{i} f_{i}\Bigr)^{D}.
    \]
\end{lemma}
\begin{proof}\leanok
    Apply the weighted arithmetic--geometric-mean inequality with uniform
    weights $1/D$:
    \[
        \Bigl(\prod_{i} f_{i}\Bigr)^{1/D} \le \frac{1}{D}\sum_{i} f_{i}.
    \]
    Raising both sides to the $D$-th power gives the stated bound.
\end{proof}

% Lean: Matrix.PosSemidef.pow_card_mul_det_re_le_trace_re_pow,
% Matrix.PosSemidef.pow_card_mul_det_re_eq_trace_re_pow_iff
\begin{lemma}[Trace--determinant AM--GM for positive-semidefinite matrices]
\label{lem:trace_det_amgm}
    \lean{Matrix.PosSemidef.pow_card_mul_det_re_le_trace_re_pow,
    Matrix.PosSemidef.pow_card_mul_det_re_eq_trace_re_pow_iff}
    \leanok
    \uses{lem:amgm_prod_sum}
    For a positive-semidefinite $D \times D$ matrix $M$,
    \[
        D^{D}\det M \le (\tr M)^{D},
    \]
    and equality holds if and only if $M = (\tr M / D)\,\mathbb{1}$.
\end{lemma}
\begin{proof}\leanok
    Both $\det M = \prod_{i}\lambda_{i}$ and $\tr M = \sum_{i}\lambda_{i}$ are
    expressed through the nonnegative eigenvalues $\lambda_{i}$ of $M$. The
    bound is the arithmetic--geometric-mean inequality with uniform weights
    $1/D$, raised to the $D$-th power; equality holds exactly when all
    eigenvalues coincide, that is, when $M$ is a scalar matrix.
\end{proof}

% Lean: Wolf.exists_normal_form_generic
\begin{theorem}[Generic normal form for CP maps]\label{thm:exists_normal_form_generic}
    \lean{Wolf.exists_normal_form_generic}
    \leanok
    \uses{def:sl_filtering, def:doubly_stochastic, lem:infimum_is_attained,
    lem:trace_det_amgm}
    Let $T : \MN{D} \to \MN{D}$ be a completely positive map whose Choi matrix
    is positive-definite (equivalently, $T$ has full Kraus rank). Then there
    exist SL-filterings $\Phi_{1}, \Phi_{2}$ such that
    $\Phi_{2} \circ T \circ \Phi_{1}$ is doubly-stochastic.

    This is the equal-dimension (square) case $T : \MN{D} \to \MN{D}$; the
    source \cite[Proposition~2.9]{Wolf2012Quantum} states the proposition for a
    general $T : \mathcal{M}_{d_1} \to \mathcal{M}_{d_2}$ with separate filterings
    on each factor.
    % See docs/paper-gaps/wolf_prop28_square_case.tex
\end{theorem}
\begin{proof}\leanok
    \uses{lem:infimum_is_attained, lem:trace_det_amgm}
    Let $\tau$ be the Choi matrix of $T$ and take the minimiser
    $(S_{1}, S_{2})$ of
    $\tr[(S_{2} \otimes_{k} S_{1})\tau(S_{2} \otimes_{k} S_{1})^{\dagger}]$ over
    $\det S_{1} = \det S_{2} = 1$ from Lemma~\ref{lem:infimum_is_attained}. Set
    $\Phi_{1}, \Phi_{2}$ to be the filterings with matrices
    $S_{1}^{\mathsf T}$ and $S_{2}$, so that the Choi matrix of
    $\Phi_{2} \circ T \circ \Phi_{1}$ equals
    $(S_{2} \otimes_{k} S_{1})\tau(S_{2} \otimes_{k} S_{1})^{\dagger}$. Holding
    one filtering factor fixed, the minimiser is optimal in the other
    coordinate, so each partial trace minimises a functional
    $\tr[S M S^{\dagger}]$ over $\det S = 1$ for a positive-semidefinite $M$,
    where
    \[
        \tr[S M S^{\dagger}] \ge D\,(\det M)^{1/D},
    \]
    with equality at the minimiser; by the equality case of
    Lemma~\ref{lem:trace_det_amgm} this forces $S M S^{\dagger} \propto \mathbb{1}$,
    hence both partial traces are proportional to the identity, which is exactly
    the doubly-stochastic condition.
\end{proof}

\begin{definition}[Diagonal Lorentz normal form]\label{def:lorentz_diagonal}
    \lean{Wolf.IsLorentzDiagonal}\leanok
    For a completely positive, trace-preserving qubit channel
    $T' : \MN{2} \to \MN{2}$, write $\widehat{T'}$ for its matrix
    in the normalized Pauli basis
    $\{\sigma_{0}/\sqrt{2},\sigma_{1}/\sqrt{2},
        \sigma_{2}/\sqrt{2},\sigma_{3}/\sqrt{2}\}$. The channel is in
    \emph{diagonal Lorentz normal form} when it is unital and every
    off-diagonal entry of $\widehat{T'}$ is zero.
\end{definition}

\begin{definition}[Non-diagonal Lorentz normal form]\label{def:lorentz_non_diagonal}
    \lean{Wolf.IsLorentzNonDiagonal}\leanok
    For a completely positive, trace-preserving qubit channel
    $T' : \MN{2} \to \MN{2}$, write $\widehat{T'}$ in the normalized
    Pauli basis. The channel is in \emph{non-diagonal Lorentz normal form}
    when, for some $x \in [0,1]$,
    \[
        \widehat{T'} =
        \begin{pmatrix}
            1   & 0          & 0          & 0   \\
            0   & x/\sqrt{3} & 0          & 0   \\
            0   & 0          & x/\sqrt{3} & 0   \\
            2/3 & 0          & 0          & 1/3
        \end{pmatrix}.
    \]
    Trace preservation supplies the first row of the displayed matrix.
\end{definition}

\begin{definition}[Singular Lorentz normal form]\label{def:lorentz_singular}
    \lean{Wolf.IsLorentzSingular}\leanok
    For a completely positive, trace-preserving qubit channel
    $T' : \MN{2} \to \MN{2}$, write $\widehat{T'}$ in the normalized
    Pauli basis. The channel is in \emph{singular Lorentz normal form} when
    \[
        \widehat{T'} =
        \begin{pmatrix}
            1 & 0 & 0 & 0 \\
            0 & 0 & 0 & 0 \\
            0 & 0 & 0 & 0 \\
            1 & 0 & 0 & 0
        \end{pmatrix},
    \]
    equivalently, every input state is mapped to $(1+\sigma_{3})/2$.
    Trace preservation supplies the first row of the displayed matrix.
\end{definition}

% Pending Lean formalization: there is currently no declaration for this theorem.
\begin{theorem}[Lorentz normal form for qubit channels]\label{thm:lorentz_normal_form_qubit}
    \uses{thm:exists_normal_form_generic,
        def:lorentz_diagonal, def:lorentz_non_diagonal, def:lorentz_singular}
    For every qubit channel $T : \MN{2} \to \MN{2}$, there exist invertible
    completely positive maps $\Phi_{1}, \Phi_{2}$, both of Kraus rank one, such
    that the filtered channel $T' = \Phi_{2} \circ T \circ \Phi_{1}$ is in one
    of the three Lorentz normal forms: diagonal, non-diagonal, or singular.
    These general filters include scalar freedom and are not restricted to
    determinant-one $\SL(2, \C)$ filterings. A proof requires the corresponding
    scalar normalization and the classification of Lorentz orbits.
\end{theorem}

%% The following sections will not be needed for the fundamental theorem of the matrix product states.

%% =========================================================
\section{Determinant of a quantum channel}
\label{sec:channel_determinant}
%% =========================================================

%% Source: \cite[Section~6.1, Theorem~6.1(2)]{Wolf2012Quantum}.

\begin{definition}[Channel determinant]\label{def:channel_det}
    \lean{channelDet}\leanok
    For a linear map $T : \MN{D} \to \MN{D}$, the
    \emph{channel determinant} $\det T$ is the determinant of the
    matrix of $T$ with respect to the standard matrix-unit basis
    of $\MN{D}$.
    This is the quantity studied in \cite[Section~6.1]{Wolf2012Quantum}.
\end{definition}

\begin{remark}[Channel determinant as an operator determinant]
    \label{rem:channel_det_linearMap_det}
    \lean{channelDet_eq_linearMap_det}
    \leanok
    The channel determinant agrees with the ordinary determinant of the linear
    endomorphism $T : \MN{D} \to \MN{D}$.
\end{remark}

\begin{definition}[Unitary channel]\label{def:unitary_channel}
    \lean{unitaryChannel, unitaryChannel_isChannel}\leanok
    For a unitary matrix $U \in \mc{U}(D)$, the
    \emph{unitary channel} is $T(\rho) = U\rho U^\dagger$.
    It is automatically a quantum channel.
\end{definition}

\begin{theorem}[Unitary channels have determinant one]
    \label{thm:channel_det_unitary_eq_one}
    \lean{channelDet_unitary_eq_one, channelDet_norm_eq_one_of_unitaryChannel}
    \leanok
    \uses{def:channel_det, def:unitary_channel}
    If $T(\rho) = U\rho U^\dagger$ is a unitary channel, then
    $\det T = 1$ and hence $|\det T| = 1$.
\end{theorem}

\begin{proof}\leanok
    Vectorization identifies $T$ with a Kronecker product
    $\overline U \otimes U$, whose determinant is
    $\overline{\det U}^{\,D}(\det U)^D = 1$.
\end{proof}

\begin{theorem}[Determinant bound for positive trace-preserving maps]
    \label{thm:channelDet_norm_le_one}
    \lean{channelDet_norm_le_one_of_positive_tracePreserving,
        channelDet_norm_le_one_of_channel}\leanok
    \uses{def:channel_det, def:positive_map, def:trace_preserving}
    For any positive trace-preserving map $T : \MN{D} \to \MN{D}$,
    $|\det T| \le 1$.
    This is \cite[Theorem~6.1(1)]{Wolf2012Quantum}.
\end{theorem}

\begin{proof}\leanok
    Every eigenvalue $\mu$ of $T$ satisfies $|\mu| \le 1$
    (positivity and trace preservation force spectral radius~$\le 1$).
    Since $\det T$ is the product of the eigenvalues
    (counted with algebraic multiplicity), the bound
    $|\det T| = \prod_i |\mu_i| \le 1$ follows by induction.
\end{proof}

\begin{theorem}[Determinant one iff the channel is unitary]
    \label{thm:channelDet_norm_eq_one_iff_unitary}
    \lean{channelDet_norm_eq_one_iff_exists_unitaryChannel}\leanok
    \uses{def:channel, def:channel_det, def:unitary_channel}
    For a CPTP map $T : \MN{D} \to \MN{D}$,
    \[
        |\det T| = 1
        \iff
        \exists\, U \in \mc{U}(D),\quad T(\rho) = U\rho U^\dagger.
    \]
    This is \cite[Theorem~6.1(2)]{Wolf2012Quantum}.
\end{theorem}

\begin{proof}
    \leanok
    \uses{thm:channelDet_norm_le_one, thm:channel_det_unitary_eq_one}
    Reverse direction is Theorem~\ref{thm:channel_det_unitary_eq_one}.
    Forward direction: $|\det T| = 1$ together with $|\mu| \le 1$ for every
    eigenvalue $\mu$ (from trace preservation and positivity) forces every
    eigenvalue to satisfy $|\mu| = 1$. Hence $T$ is bijective. A bijective
    CPTP map is a $*$-automorphism: the Heisenberg dual
    $T^*(Y) = \sum_i K_i^\dagger Y K_i$ is unital (since $T$ is TP), so the
    Kadison--Schwarz inequality
    $T^*(Y^\dagger Y) \ge T^*(Y)^\dagger T^*(Y)$ becomes an equality for every
    $Y$. Hence $T^*$, and therefore $T$, preserves products. By the
    Skolem--Noether theorem, every $*$-automorphism of $\MN{D}$ has the form
    $T(\rho) = U \rho U^\dagger$ for some unitary $U$. Kraus freedom then
    gives $K_i = c_i U$ with $\sum_i |c_i|^2 = 1$.
\end{proof}
